\documentclass[10pt,a4paper]{amsart}
\usepackage[margin=19mm]{geometry}
\usepackage{amsmath,amssymb,mathtools}
\usepackage[scr=rsfs,cal=euler]{mathalfa}
\usepackage{xfrac}
\usepackage[unicode,hidelinks,bookmarks=false]{hyperref}
\newcommand{\ie}{i.e.\ }
\newcommand{\cf}{cf.\ }
\newcommand{\resp}{respectively\ }
\newcommand{\Subsection}{\subsection}
\providecommand{\nobreakdash}{\nobreak\hbox{-}}
\setlength{\emergencystretch}{2em}
\multlinegap0pt
\usepackage{mathtools,amscd}
\usepackage{mathrsfs}
\usepackage{color}
\usepackage{tikz}
\usepackage{float}
\usepackage{cleveref}
\usetikzlibrary{positioning}
\newcommand{\field}[1]{\mathbb{#1}}
\newcommand{\Z}{\field{Z}}
\newcommand{\Q}{\field{Q}}
\newcommand{\R}{\field{R}}
\newcommand{\C}{\field{C}}
\newcommand{\PGL}{\mathrm{PGL}}
\newcommand{\PSL}{\mathrm{PSL}}
\newcommand{\F}{\field{F}}
\newcommand{\Oint}{\mathcal{O}}
\newcommand{\ip}{\mathrm{i}_{p}}
\newcommand{\iK}{\mathrm{i}_{K}}
\newcommand{\iphi}{\mathrm{i}_{\varphi}}
\newcommand{\Gal}{\mathrm{Gal}}
\renewcommand{\vec}[1]{\boldsymbol{#1}}
\newtheorem{theorem}{Theorem}[section]
\newtheorem{lemma}[theorem]{Lemma}
\newtheorem{proposition}[theorem]{Proposition}
\newtheorem{corollary}[theorem]{Corollary}
\newtheorem{algorithm}[theorem]{Algorithm}
\newtheorem{example}[theorem]{Example}
\newtheorem{definition}[theorem]{Definition}
\newtheorem{problem}[theorem]{Problem}
\newtheorem{property}[theorem]{Property}
\newtheorem{remark}[theorem]{Remark}
\newtheorem*{notation}{Notation}
\def\labelenumi{\theenumi}
\def\theenumi{(\roman{enumi})}
\def\labelenumii{\theenumii}
\def\theenumii{(\arabic{enumii})}
\def\labelenumiii{\theenumiii}
\def\theenumiii{(\alph{enumiii})}
\DeclareMathOperator{\gal}{Gal}
\allowdisplaybreaks
\numberwithin{equation}{section}
  \def\gn#1#2{{${http://groupnames.org/\#?#1}{#2}$}}
\def\gn#1#2{$#2$}  % comment this line out to get html links
\tikzset{sgplattice/.style={inner sep=1pt,norm/.style={red!50!blue},char/.style={blue!50!black},
  lin/.style={black!50}},cnj/.style={black!50,yshift=-2.5pt,left=-1pt of #1,scale=0.5,fill=white}}
\begin{document}
\title[Computating decomposition groups and inertia groups using Ne]{Computating decomposition groups and inertia groups using Newton polygons}
\author{Kazuma Igarashi; Nozomu Suzuki}
\date{}
\maketitle
\noindent\emph{Source and licence.} Computating decomposition groups and inertia groups using Newton polygons. Version arXiv:2606.26888v1 (CC BY 4.0). This material is distributed under the Creative Commons Attribution 4.0 International licence, http://creativecommons.org/licenses/by/4.0/, as recorded in the arXiv OAI licence metadata for this exact version and shown on the abstract page https://arxiv.org/abs/2606.26888v1. Full rights notice: licenses/arxiv-2606.26888-CC-BY-4.0.txt.\par\smallskip
\noindent\textit{Editorial scope.} Counted unit: the original \texttt{theorem} environment \texttt{thm:our main theorem} at source lines 502--518 together with its complete proof at lines 802--874; the delivered document keeps source lines 125--875 so that every referenced label and every citation resolves inside it. Native editable source is the paper's own concatenated TeX; the AMS wrapper is editorial formatting only. No publisher class, upstream script or unrelated artwork is redistributed.
\section{Preliminaries}\label{sec:Prelim}

  Let $p$ be a prime number. Fix an algebraic closure $\overline{\Q}_p$ (resp. $\overline{\field{F}}_p$) of $\field{Q}_p$ (resp. $\field{F}_p$). 
  Throughout, every algebraic extension of $\field{Q}_p$ (resp. $\field{F}_p$) is 
  regarded as a subfield of $\overline{\field{Q}}_p$ (resp. $\overline{\field{F}}_p$). 
  Let $K$ be a finite extension of $\field{Q}_p$. We denote by $\mathcal{O}_K$ the ring of integers of $K$, 
  by $\mathfrak{p}_K$ the prime ideal of $\mathcal{O}_K$, by $\field{F}_K$ the residue field, and by $v_K$ the valuation of $\overline{\Q}_p$ 
  normalized by $v_K(K^{\times}) = \Z$. 
  We fix a uniformizer $\pi \in \mathcal{O}_K$. 
  For $a \in \mathcal{O}_K$, the image of $a$ in $\field{F}_K$ is denoted by $\overline{a}$. 
  We shall denote the image of $\varphi(X) \in \mathcal{O}_K [X]$ in $\field{F}_K [X]$ also by $\overline{\varphi}(X)$. 
  In this paper, we assume that a polynomial $f$ is separable.

  \subsection{Newton polygons}\label{sec:Newton polygons}
    
    We define $\varphi$-polygons, which generalize the classical Newton polygon, and the polynomials associated to them. 
    To this end, we extend the valuation $v_K$ to $\mathcal{O}_K[X]$, which we denote by the same symbol, by 
    \begin{align}
      v_K:\mathcal{O}_K[X] & \rightarrow \Z_{\ge 0} \cup \{ \infty \}, \\
      a_0 + a_1 X + \cdots + a_n X^n & \mapsto \min\{v_K(a_i):0 \le i \le n \}. 
    \end{align}
    Let $\varphi(X) \in \mathcal{O}_K[X]$ be a monic irreducible polynomial of degree $l \ge 1$ such that 
    its reduction $\overline{\varphi}(X)$ is also irreducible over $\field{F}_K$. 
    Fix a root $\zeta$ of $\overline{\varphi}$. 
    Given a polynomial $f(X) \in \mathcal{O}_K[X]$ of degree $n$, we obtain a unique \emph{$\varphi$-expansion} as 
    \[
      f(X) = \sum_{i = 0}^{[n/l]} a_i(X) \varphi(X)^i, 
    \]
    where each $a_i(X) \in \mathcal{O}_K[X]$ is a polynomial of degree $< l$ or $a_i(X)=0$. 
    Throughout, we assume $a_0(X) \ne 0$. 

    When $f(X) \in \mathcal{O}_K[X]$ is expanded as above, we define the \emph{$\varphi$-polygon} of $f(X)$ 
    as the lower convex envelope in the Euclidean plane of points $\{(i,v_K(a_i))|0 \le i \le [n/l]\}$. 
    When $\varphi(X) = X$, the $\varphi$-polygon is the classical Newton polygon. 
    If $f(X)$ is a monic polynomial, we call the set of sides with negative slope the \emph{principal part} of $\varphi$-polygon. 
    When $\overline{\varphi}$ does not divide $\overline{f}$, the $\varphi$-polygon of $f$ has no principal part. 
    We say that the $\varphi$-polygon is \emph{one-sided} when it has exactly one side; otherwise, we say that it is \emph{many-sided}. 

    In general, let $S$  be a line segment in $\R^2$ with negative slope and its endpoints
    having integer coordinates. 
    Let its initial point be $(r,s)$ and its terminal point be $(r + E,s - H)$. 
    We call $E$ the \emph{length} of $S$ and $H$ the \emph{height} of $S$, denoted $E(S)$ and $H(S)$, respectively. 
    We define 
    \[
      d = \gcd(E(S),H(S)), 
    \]
    and call $d$ the degree of $S$. If $H = E = 0$, we set $d = 0$. 

    Let $f(X) \in \mathcal{O}_K[X]$ be a monic polynomial, and $S$ be a line segment 
    in $\{ (x,y) \in \R^2 \mid x \ge 0, y \ge 0 \}$ with negative slope and its endpoints having integer coordinates. 
    Suppose that no vertex of the $\varphi$-polygon of $f$ lies below $S$. 
    We denote the endpoints of $S$ by $P = (e_0,h_0)$ and $Q = (e_1,h_1)$, where $e_0 \le e_1$. 
    Write $d$ for the degree of $S$. 
    Put 
    \[
      b_j(X) = \pi^{-h_0 + jh} a_{e_0 + je}, 
    \]
    where $h = \frac{h_0 - h_1}{d}, e = \frac{e_1 - e_0}{d}$. 
    We define the \emph{polynomial associated to $f$ and $S$}, or the \emph{associated polynomial of $f$ and $S$} 
    as 
    \[
      f_S(Y) = \sum_{j=0}^d \overline{b_j}(\zeta) Y^j.
    \]
    If no vertex lies on $S$, we set $f_S(Y) = 0$. 
    We also define the normalized associated polynomial by 
    \[
      f_S^{\mathrm{norm}}(Y) = \overline{b_d}(\zeta)^{-1} \sum_{j=0}^d \overline{b_j}(\zeta) Y^j.
    \]

    We now define the associated polynomial for a general line or line segment $S$. 
    Let $V$ be the set of integer lattice points lying on $S \cap \{ (x,y) \in \R^2 \mid x \ge 0, y \ge 0 \}$. 
    We denote by $P, Q$ the points of $V$ with minimal and maximal abscissa, respectively, and 
    set $S'$ to be the line segment between $P$ and $Q$. 
    We define the associated polynomials of $f$ and $S$ by $f_S = f_{S'}$ and $f^{\mathrm{norm}}_S = f^{\mathrm{norm}}_{S'}$. 

    \begin{remark}
      If $f = \varphi^m$ and the $\varphi$-polygon of $f$ is one-sided, then $f_S = f_S^{\mathrm{norm}}$. 
    \end{remark}

  \subsection{Works of Ore, Montes and Nart on Newton polygons}\label{sec:Ore Montes Nart}
    
    In this section, we review the works of Ore and of Montes--Nart on Newton polygons. 

    The following two theorems are well-known facts about Newton polygons.

    \begin{theorem}[Theorem of the product]\label{thm:theorem of product}
      Let $f_1, \dots ,f_g \in \mathcal{O}_K[X]$ be monic polynomials, and write 
      $f(X) = f_1(X) \cdots f_g(X)$. 
      Let $S_{i,1}, \dots S_{i,k_i}$ be the sides of the principal part of the $\varphi$-polygon 
      of $f_i$. We define $M=\{-H(S_{i,j})/E(S_{i,j}) \in \Q_{< 0} \mid 1 \le i \le g,1 \le j \le k_i\}$, 
      and let $M'$ be the set of slopes of the sides in the principal part of the $\varphi$-polygon of $f$. 
      Then, $M=M'$, and for any side $S$ of the $\varphi$-polygon of $f$ with slope $-m$, we have 
      \[
        E(S)=\sum_{\substack{i,j\\H(S_{i,j})/E(S_{i,j})=m}} E(S_{i,j}), H(S)=\sum_{\substack{i,j\\H(S_{i,j})/E(S_{i,j})=m}} H(S_{i,j}). 
      \]
      Moreover, for any $m \in \Q_{>0}$, let $S_m$ be a side of the $\varphi$-polygon of $f$ of slope $-m$, 
      and let $S_{i,m}$ be a side of the $\varphi$-polygon of $f_i$ of slope $-m$. 
      Then, 
      \[
        f_S^{\mathrm{norm}}(Y) = \sideset{}{'}\prod_{i}f_{S_{i,m}}^{\mathrm{norm}}(Y), 
      \]
      where the product runs over all $i$ such that the $\varphi$-polygon of $f_i$ has a side of slope $-m$. 
    \end{theorem}

    \begin{theorem}[Theorem of the polygon]\label{thm:theorem of the polygon}
      Let $f \in \mathcal{O}_K[X]$ be a monic polynomial whose reduction is a power of $\varphi$ and whose $\varphi$-polygon
      consists of the sides $S_1, \dots ,S_g$. Let the slope of $S_{i}$ be $-m_i$. 
      Then, there exist monic polynomials $f_1, \dots ,f_g \in \mathcal{O}_K[X]$ 
      satisfying the following conditions: 
      \begin{enumerate}
        \item $ f(X) = \displaystyle \prod_{i = 1}^g f_i(X) $. 

        \item The $\varphi$-polygon of $f_i$ is one-sided, and it has the same length and height as $S_{i}$. 

        \item Let $S_{i}$ also denote the $\varphi$-polygon of $f_i$, then $f_{S_{i}}^{\mathrm{norm}}(Y) = (f_i)_{S_{i}}^{\mathrm{norm}}(Y)$. 

        \item If $\theta \in \overline{\Q}_p$ is a root of $f_i(X)$, then $v_K(\varphi(\theta)) = m_i$. 
      \end{enumerate}
    \end{theorem}

    Let $f \in \mathcal{O}_K[X]$ be a monic polynomial whose $\varphi$-polygon is one-sided, and denote this polygon by $S$. 
    We say $f$ is \emph{regular} if the associated polynomial $f_S(Y)$ is separable. 
    For a regular polynomial, Ore showed that the ramification index and inertia degree of the field defined by the polynomial can be computed from the $\varphi$-polygon. 

    \begin{theorem}[Theorem of Ore]\label{thm:theorem of Ore}
      Let $f(X) \in \mathcal{O}_K[X]$ be a monic polynomial whose $\varphi$-polygon is one-sided, and denote this polygon by $S$. 
      Suppose that $\overline{f}(Y) = \overline{\varphi}(Y)^n (n \in \Z_{\ge 0})$. 
      We write $e=E(S)/d$ and $h=H(S)/d$. Let 
      \[
        f_S(Y) = \psi_1(Y)^{e_1} \cdots \psi_t(Y)^{e_t} 
      \]
      be the factorization of $f_S^{\mathrm{norm}}$ into a product of powers of distinct irreducible polynomials
      of $\field{F}_K[Y]$. 
      Then $f$ admits a factorization 
      \[
        f(X) = f_1(X) \cdots f_t(X), 
      \]
      where each $f_i(X)$ is a monic polynomial with a one-sided $\varphi$-polygon $S_{i}$ 
      of the same slope as $S$, and whose associated polynomial $(f_i)_{S_{i}}(Y)$ equals $\psi_i(Y)^{e_i}$. 
      Moreover, we suppose that $f(X)$ is regular (i.e., $e_1 = \cdots = e_t = 1$), and let $\theta$ be a root of $f_i(X)$ and $L = K(\theta)$. 
      Then, all $f_i$ are irreducible, the prime ideal of $L$ equals $(\varphi(\theta)^b/\pi^c)\mathcal{O}_L$ 
      where $b$ and $c$ are positive integers such that $bh - ce = 1$, and $\mathcal{O}_L$ is the ring of integers of $L$, and 
      \[
        e(L/K) = e, f(L/K) = l \cdot \deg \psi_i(Y). 
      \]
    \end{theorem}

    Now, we define the two indices $i_K(f)$ and $i_\varphi(f)$ of a polynomial. 

    \begin{definition}\label{def:index of K}
      For a monic polynomial $f(X) \in \mathcal{O}_K[X]$, we define $i_K(f)$ as follows: 
      \begin{itemize}
        \item If $f(X)$ is irreducible, choose a root $\theta$, and put $L=K(\theta)$. 
          We then define 
          \[
            i_K(f) = v_K([\mathcal{O}_L:\mathcal{O}_K[\theta]])/[K:\Q_p]. 
          \]
        \item For monic irreducible polynomials $g(X),h(X) \in \mathcal{O}_K[X]$, 
          we set 
          \[
             r(g,h) = v_K(\mathrm{Res}(g,h)), 
          \]
          where $\mathrm{Res}(g,h)$ is the resultant of $g(X)$ and $h(X)$. 
        \item Let $f(X)$ be a product of monic irreducible polynomials $f_1(X), \dots ,f_t(X) \in \mathcal{O}_K[X]$. 
          We define 
          \[
            i_K(f) = \sum_{i = 1}^t i_K(f_i) + \sum_{i<j} r(f_i,f_j). 
          \]
      \end{itemize}
    \end{definition}

    \begin{definition}\label{def:index of phi}
      Let $f(X) \in \mathcal{O}_K[X]$ be a monic polynomial, and suppose that the principal part 
      of the $\varphi$-polygon of $f(X)$ consists of the sides $S_1, \dots ,S_r$. 
      For each $i$, denote the length, height, and degree of $S_{i}$ by $E_i$, $H_i$, and $d_i$, respectively. 
      Then, we define 
      \[
        i_{\varphi}(f) = \deg\varphi\cdot\left(\sum_{i = 1}^{r - 1} E_i \cdot 
        \left(\sum_{j = i + 1}^r H_j\right) + \frac{1}{2} \sum_{i = 1}^r 
        (H_i E_i - H_i - E_i + d_i)\right).
      \]
    \end{definition}

    The quantity $i_{\varphi}(f)/\deg\varphi(X)$ represents the number of points of 
    integer coordinates below or on the $\varphi$-polygon of $f$, excluding those lying on both axes. 
    See \cref{fig:1}.

    \begin{figure}
      \centering
      \begin{tikzpicture}
        \draw[->] (0,-0.5) -- (0,3);
        \draw[->] (-0.5,0) -- (5,0);
        \draw (0,2.5) -- (1,1.5) -- (3,0.5) -- (4.5,0);
        \foreach \x in {0.5,1,1.5,2}
          \fill[black] (0.5,\x) circle [radius=0.05];
        \foreach \x in {0.5,1,1.5}
          \fill[black] (1,\x) circle [radius=0.05];
        \foreach \x in {0.5,1}
          \fill[black] (1.5,\x) circle [radius=0.05];
        \foreach \x in {0.5,1}
          \fill[black] (2,\x) circle [radius=0.05];
        \foreach \x in {0.5}
          \fill[black] (2.5,\x) circle [radius=0.05];
        \foreach \x in {0.5}
          \fill[black] (3,\x) circle [radius=0.05];
      \end{tikzpicture}
      \caption{The number of points of integer coordinates}
      \label{fig:1}
    \end{figure}

    Let $f(X) \in \mathcal{O}_K[X]$ be a monic irreducible polynomial. 
    By Hensel's lemma, there exists a monic irreducible polynomial $\overline{\varphi}(Y) \in \field{F}_K[Y]$ 
    such that $\overline{f}(Y) = \overline{\varphi}(Y)^n (n \in \Z_{> 0})$. 
    Assume that $f(X)$ has a one-sided $\varphi$-polygon $S$, and set $e=E(S)/d$, $h=H(S)/d$. 
    By Theorem \ref{thm:theorem of Ore}, there exists a monic irreducible polynomial $\psi(Y) \in \field{F}_K[Y]$ 
    such that $f_S(Y) = \psi(Y)^a (a \in \Z_{> 0})$. Let $\theta$ be a root of $f(X)$ and put $L=K(\theta)$. 
    We define $\gamma_i=\varphi(\theta)^i/\pi^{[iH/E]} \in \mathcal{O}_L \quad (i = 0, \dots ,n - 1)$. 

    \begin{proposition}\label{prop:prop of Montes Nart}\cite[Proposition]{MR1152908}
      Under the above assumptions, the following conditions are equivalent: 
      \begin{enumerate}
        \item $i_K(f) = i_{\varphi}(f).$

        \item $e(L/K) = ea$, $f(L/K) = l \cdot \deg\psi(Y)$ and 
        [$a = 1$ or $\mathfrak{p}_L = \varphi(\gamma_e)\mathcal{O}_L$]. 
      \end{enumerate}
    \end{proposition}

    To prove the above proposition, Montes and Nart proved the following lemma on the valuation of the resultant. 
    
    \begin{lemma}\label{lem:valuation of resultant}\cite[Lemma 2]{MR1152908}
      Let $f,g \in \mathcal{O}_K[X]$ be monic polynomials of degrees $n$, $n'$, respectively. 
      Assume that the $\varphi$-polygons $S$ and $S'$ of $f$ and $g$ are one-sided, and denote their heights by $H$ and $H'$, respectively. 
      Then, 
      \[
        v_K(\mathrm{Res}(f,g)) \ge \min\{nH',n'H\}. 
      \]
      Moreover, equality holds if and only if either slope of $S$, $S'$ are different or $\mathrm{Res}(f_S,g_{S'}) \ne 0$. 
    \end{lemma}

     Using \cref{prop:prop of Montes Nart}, Montes and Nart proved the following theorem, which is a generalization of theorem of Ore.

    \begin{theorem}\label{thm:theorem of Montes Nart}\cite[Theorem 1]{MR1152908}
      Under the assumptions and notation of Theorem \ref{thm:theorem of Ore}, suppose that $i_K(f) = i_{\varphi}(f)$. 
      Then, $f_1(X), \dots ,f_t(X)$ are irreducible, and 
      \[
        e(L/K) = ee_i, f(L/K) = l\cdot\deg\psi_i(Y). 
      \]
    \end{theorem}

    We can check the condition $i_K(f) = i_{\varphi}(f)$ as follows.

    We continue to use the above notation. 
    Let $T$ be the unique unramified extension of $K$ of degree $l$. 
    Let $\mathcal{O}_K[X,Y] \to \field{F}_K[X,Y]/(\bar{\varphi}(X)) \cong \field{F}_T[Y]$ be the natural homomorphism, 
    and choose a lift $\Psi_i(X,Y)$ of $\psi_i(Y)$ such that 
    \[
      \Psi_i(X,Y) = \sum_j a_j(X)Y^j,\quad \deg a_j(X) < \deg\varphi(X) \quad \text{or} \quad a_j(X) = 0. 
    \]
    We define 
    \[
      f^0(X) = \pi^H \prod_{i = 1}^t \Psi_i(X,\varphi(X)^e/\pi^h)^{e_i} \in \mathcal{O}_K[X], 
    \]
    and 
    \[
      f^1(X) = f(X) - f^0(X). 
    \]
    Let $\tilde{S}$ be the side $S$ shifted upwards by $1/e$. 
    Then, the $\varphi$-polygon of $f^1(X)$ has no points below $\tilde{S}$. 

    \begin{theorem}\label{thm:equivalent condition of Montes Nart}\cite[Criterion 1]{MR1152908}
      The equality $i_K(f) = i_{\varphi}(f)$ holds if and only if we have either 
      $a_i = 1$ or $\psi_i(Y) \nmid f^1_{\tilde{S}}(Y)$ in $\F_T[Y]$ for all $i$. 
    \end{theorem}

    Let $f(X) \in \mathcal{O}_K[X]$ be a monic polynomial, not necessarily one-sided.
    Let $S_1, \dots ,S_r$ be the sides of the principal part of the $\varphi$-polygon of $f$. 
    By \cref{thm:theorem of the polygon}, $f(X)$ admits a factorization 
    $f(X)=f_1(X) \cdots f_r(X)$ where the $\varphi$-polygon of $f_i$ is $S_{i}$. 
    Since each $f_i$ is one-sided, we can define $f_i^0$ and $f_i^1$ as above. 
    We then define 
    \[
      f^1(X)=f(X)-LC_{\varphi}(f) \cdot \prod_{i=1}^r f_i^0(X), 
    \]
    where $LC_{\varphi}(f)$ is the leading coefficient of the $\varphi$-expansion of $f$. 
    Now, we define $\widetilde{S_{i}}$ again.
    Let $\overline{S}_{i}$ be a line containing $S_{i}$.
    Put $\widetilde{S_{i}}$ be the minimal line segment containing $\overline{S}_{i} \cap (\Z_{\ge 0} \times \Z_{\ge 0})$. 

    \begin{theorem}\label{thm:criterion 2 of Montes Nart}\cite[Criterion 2]{MR1152908}
      We use the above notation. Let $\psi(Y)$ be an irreducible factor of $(f_i)_{S_{i}}$. 
      Then, $\psi(Y)$ divides $(f_i^1)_{\widetilde{S_{i}}}(Y)$ if and only if it divides $(f^1)_{\widetilde{S_{i}}}(Y)$. 
    \end{theorem}

    To apply \cref{thm:theorem of Montes Nart} to the many-sided case, 
    we have to check the condition on indices. 

    \begin{definition}
        Let $f(X) \in \Oint_{K}[X]$.
        We call a set $\{\varphi_{i}(X)\}_{i \in I}$ of monic irreducible polynomials in $\Oint_{K}[X]$
        a \emph{set of irreducible pullbacks} of $f$
        if $\{\overline{\varphi_{i}}(X)\}_{i \in I}$ coincides with 
        the set of the distinct irreducible factors of $\overline{f}(X) \in \F_{K}[X]$.
    \end{definition} 

    The following theorem is an analogue of \cref{thm:equivalent condition of Montes Nart}. 

    \begin{theorem}[{\cite[Theorem 2]{MR1152908} \cite[Theorem 3.1]{index}}]\label{thm:theorem of index}
        Let $\Phi$ be a set of irreducible pullbacks of $f$.
        Then,
        \[
            \iK(f) \ge \sum_{\varphi \in \Phi} \iphi(f). \label{1:ineq:thm}
        \]
        Moreover, the equality holds if and only if for any $\varphi(X) \in \Phi$
        we have either $\Psi(Y)^2 \nmid f_{S}(Y)$ or $\Psi(Y) \nmid f^{1}_{\tilde{S}}(Y)$
        for any side $S$ of the principal parts
        of the $\varphi$-polygons and
        for any irreducible factor $\Psi(Y)$ of $f_{S}(Y)$.
    \end{theorem}

  \subsection{Works of Kölle and Schmid}\label{sec:Kölle Schmid}

    Now, we review the theorem of Kölle and Schmid. 
    Throughout this section, we suppose that $\overline{f}(X) = X^n$ and $\varphi(X) = X$. 

    In this section, we assume that $K$ is a (global) number field. 
    Let $f(X)=\sum_{i = 0}^n a_i X^{n - i} \in \mathcal{O}_K[X]$ be a polynomial, $L$ be the splitting field of $f(X)$, and $G = \Gal_K(f) = \Gal(L/K)$. 
    Let $\mathfrak{p}$ be a prime ideal of $K$. 
    We denote the set of all roots of $f$ by $Z_f$. Then $G$ acts on $Z_f$. 
    Let $\mathfrak{P}$ be a prime ideal of $L$ lying above $\mathfrak{p}$, and $G_\mathfrak{P}$ and $I_\mathfrak{P}$ be the 
    decomposition group and the inertia group of $L/K$ at $\mathfrak{P}$, respectively. 
    Then, $G_{\mathfrak{P}} = \Gal(L_{\mathfrak{P}}/K_{\mathfrak{p}})$. 
    Denote the set of roots of $f$ whose $\mathfrak{p}$-adic valuation is equal to $-m$ by $Z_{f,m}$. 
    Then $G_{\mathfrak{P}}$ acts on $Z_{f,m}$. 
    We set $G_{\mathfrak{P}}^{Z_{f,m}} = G_{\mathfrak{P}}/\mathrm{Stab}_{G_{\mathfrak{P}}}(Z_{f,m})$. 
    $G_{\mathfrak{P}}^{Z_{f,m}}$ also acts on $Z_{f,m}$. 
    Let $S_m$ be the side of $\varphi$-polygon $S$ of $f(X)$ with slope $m$. 
    Put $e=E(S_m)/d$ and $h=H(S_m)/d$, and let $(s,v_{\mathfrak{p}}(a_s))$ be its initial point. 
    We define $f_m =  \displaystyle{\sum_{(i,v_{\mathfrak{p}}(a_i)) \in S_m}}  (-1)^i a_s^{-1} a_{i}X^{de-i}$. 
    Let the distinct irreducible factors of $f_S^{\mathrm{norm}}$ over $\F_{K_{\mathfrak{p}}}$ have 
    degrees $d_1, \dots ,d_r$ (so that $\sum_{i = 1}^r d_i = d = \deg (f_S^{\mathrm{norm}})$). 
    
    The following is the theorem of Kölle and Schmid. 

    \begin{theorem}[{\cite[Theorem]{MR2102807},\cite{MR2557856}}]\label{thm:theorem of Kölle Schmid}
      We use the above notation, and suppose that $f$ is regular. 
      Then, $f_m$ is separable and $|Z_{f,m}| = de$. 
      If in addition $\mathfrak{p} \nmid e$, then the following hold:
      \begin{enumerate}
        \item For every root $\beta$ of $f_m$, there exists $\theta \in Z_{f,m}$ such that $K_{\mathfrak{p}}(\beta) = K_{\mathfrak{p}}(\theta)$, and vice versa. 
        \item $G_{\mathfrak{P}}^{Z_{f,m}}$ is permutation isomorphic to $\Gal_{K_\mathfrak{p}}(f_m)$. 
        \item $I_{\mathfrak{P}}^{Z_{f,m}}$ is cyclic and generated by an element which is the product of $d$ disjoint $e$-cycles on $Z_{f,m}$. 
        \item $G_{\mathfrak{P}}^{Z_{f,m}} = \langle \sigma,\tau \rangle$ has just $r$ orbits of sizes $d_1 e, \dots ,d_r e$ 
          and is a metacyclic group, with $\sigma^{-1}\tau\sigma = \tau^{N\mathfrak{p}}$. 
          Let $\mu = \mathrm{lcm}(\omega,d_1, \dots ,d_r) \quad (\omega = \mathrm{ord}_e N\mathfrak{p})$, the order of the (cyclic) group $G_{\mathfrak{P}}^{Z_{f,m}}/I_{\mathfrak{P}}^{Z_{f,m}}$ 
          is divisible by $\mu$, and it is a divisor of $\mu e$. 
          This order is equal to $\mu$ if $e = 1$ and $d = 1$, and if $r = 1$ and $\gcd (\omega,d) = 1$. 
      \end{enumerate}
    \end{theorem}


\section{Main theorem}\label{sec:Main theorem}

  \subsection{One-sided case}\label{sec:One-sided}

  In this section, we prove a generalization of the theorem of Kölle and Schmid for polynomials having a one-sided $\varphi$-polygon. 
  Throughout this section, we assume that $\overline{f}(X) = X^n$ and $\varphi(X) = X$. 

  To state the theorem, we define the following polynomial. 
  We use the notation of \cref{sec:Prelim}. 
  Let $f(X) \in \mathcal{O}_K[X]$ be a monic polynomial having a one-sided polygon $S$. 
  We define $\tilde{f}^1(X) \in \mathcal{O}_K[X]$ by $\tilde{f}^1(X) = (f^1)^0(X)$, 
  using the definition of \cref{sec:Ore Montes Nart} applied to $f^1$ and its $\varphi$-polygon $\tilde{S}$. 
  We put $\tilde{f}=f^0 + \tilde{f}^1$. 

  In the following theorem, we use the notation of \cref{sec:Kölle Schmid}. 


  \begin{theorem}\label{thm:our main theorem}
      Suppose that $f$ has a one-sided $\varphi$-polygon and $i_{K_{\mathfrak{p}}}(f) = i_{\varphi}(f)$. 
      Then, $\tilde{f}$ is separable, $f$ has $de$ roots, and the $\mathfrak{p}$-adic valuations of all roots of $f$ are equal to $m$. 
      Moreover, let $f_S = \psi_1^{a_1} \cdots \psi_r^{a_r}$ be the irreducible factorization over $\field{F}_{K_{\mathfrak{p}}}$, 
      and put $a = \mathrm{lcm}\{a_1, \dots ,a_r\}$. 
      If in addition $\mathfrak{p} \nmid ea$, then the following hold:
      \begin{enumerate}
        \item For every root $\beta$ of $\tilde{f}$, there exists $\theta \in Z_f$ such that $K_{\mathfrak{p}}(\beta) = K_{\mathfrak{p}}(\theta)$, and vice versa. 
        \item $G_{\mathfrak{P}}$ is permutation isomorphic to $\Gal_{K_\mathfrak{p}}(\tilde{f})$. 
        \item $I_{\mathfrak{P}}$ is cyclic and generated by an element which is the product of $d$ disjoint $ea$-cycles on $Z_f$. 
        \item $G_{\mathfrak{P}} = \langle \sigma,\tau \rangle$ has just $r$ orbits of sizes $d_1 ea_1, \dots ,d_r ea_r$ 
          and is a metacyclic group, with $\sigma^{-1}\tau\sigma = \tau^{N\mathfrak{p}}$. 
          Let $\mu = \mathrm{lcm}(\omega,d_1, \dots ,d_r) \quad (\omega = \mathrm{ord}_{ea} N\mathfrak{p})$, the order of the (cyclic) group $G_{\mathfrak{P}}/I_{\mathfrak{P}}$ 
          is divisible by $\mu$, and it is a divisor of $\mu ea$. 
          % This order is equal to $\mu$ if $ea = 1$ or $d = 1$, and if $r = 1$ and $\gcd (\omega,d) = 1$. 
      \end{enumerate}
    \end{theorem}


  To prove this theorem, we prove several propositions on polynomials having a one-sided $\varphi$-polygon. 
  Throughout this section, we assume that $i_K(f) = i_{\varphi}(f)$. 
  Since all points of $\varphi$-polygons of $f(X)$ and $\tilde{f}(X)$ lying on $S$ coincide, 
  we have $f_S(Y) = \tilde{f}_S(Y)$. 

  \begin{proposition}\label{prop:claim 2}
    $\tilde{f}(X)$ is separable. 
  \end{proposition}

  \begin{proof}
    Note that $f_S(Y) = \tilde{f}_S(Y)$ and ${f^1}_{\tilde{S}}(Y) = {\tilde{f}^1}_{\tilde{S}}(Y)$. 
    By \cref{thm:theorem of index}, 
    \[
      i_K(f) = i_{\varphi}(f) \Leftrightarrow i_K(\tilde{f}) = i_{\varphi}(\tilde{f}). 
    \]
    Therefore, by \cref{thm:theorem of Montes Nart}, $\tilde{f}(X)$ factors into distinct irreducible polynomials. 
  \end{proof}

  Let $F_f$ and $F_{\tilde{f}}$ be the sets of irreducible factors of $f$ and $\tilde{f}$, respectively. 
  The following proposition follows from \cref{thm:theorem of Montes Nart}. 

  \begin{proposition}\label{prop:claim 3}
    There exists a one-to-one correspondence $\phi \leftrightarrow \psi$ between $F_f$ and $F_{\tilde{f}}$ satisfying the following conditions: 
    \begin{enumerate}
      \item There exists an irreducible polynomial $g(Y) \in \field{F}_K[Y]$ such that $\phi_{S_1}(Y) = \psi_{S_2}(Y) = g(Y)^a$, where $S_1$ and $S_2$ are the $\varphi$-polygons of $\phi$ and $\psi$, respectively. 
      \item Let $\theta$ and $\beta$ be roots of $\phi$ and $\psi$, respectively. Then, 
      \begin{align}
        e(K(\theta)/K) &= e(K(\beta)/K) = ea, \\
        f(K(\theta)/K) &= f(K(\beta)/K) = \deg(g). 
      \end{align}
    \end{enumerate}
  \end{proposition}

  \begin{proposition}\label{prop:claim 4}
    Suppose that $f_S = \psi^a$ for some irreducible polynomial $\psi \in \field{F}_K[Y]$, and 
    that $\mathfrak{p} \nmid ea$. 
    Then there exists a one-to-one correspondence $\theta \leftrightarrow \beta$ between the roots of $f$ and those of $\tilde{f}$ such that $K(\theta) = K(\beta)$. 
  \end{proposition}

  To prove this proposition, we prove the following lemma. 

  \begin{lemma}\label{lem:classification of roots}
    Let $\{\alpha_1, \dots ,\alpha_t\}$ be the set of roots of $f_S(Y)$ in $\overline{\field{F}}_K$ 
    with multiplicities $a_1, \dots ,a_t$, respectively. Assume that $\mathfrak{p} \nmid e$. 
    Let $Z_f$ be the set of roots of $f$ in $\overline{K}$. 
    We define an equivalence relation on $Z_f$ by $v_K(\beta - \beta') > \frac{h}{e} \quad (\beta, \beta' \in Z_f)$ and 
    denote the decomposition of $Z_f$ into equivalence classes by $Z_f = \coprod_{i = 1}^l Z_f^i$. 
    Then $l = et$, and $\#Z_f^i = a_i$. 
    In particular, if $f_S = \psi^a$ for some irreducible polynomial $\psi \in \field{F}_K[Y]$, 
    we have $\#Z_f^i = a$ for all $i$. 
  \end{lemma}

  \begin{proof}
    We first show that it is sufficient to consider the case that $f_S$ is a power of a polynomial of degree one. 
    Let $\hat{T}$ be the unique unramified extension of $K$ of degree $\deg f_S$, then 
    $f_S = \prod_{k = 1}^t {\psi_k}^{a_k}$ for some distinct linear polynomials $\psi_k \in \field{F}_{\hat{T}}[Y] (k = 1, \dots ,t)$. 
    By \cref{thm:theorem of Ore}, $f$ admits a factorization $f = \prod_{k = 1}^t f_k$ over $\hat{T}$ such that $(f_k)_{S_{k}} = {\psi_k}^{a_k}$, where $S_{k}$ is the $\varphi$-polygon of $f_k$. 
    Denote the decomposition of $Z_{f_k}$ into equivalence classes by $Z_{f_k} = \coprod_{i = 1}^{l_k} Z_{f_k}^i$. 
    We will show that $Z_{f_k}^i$ forms an equivalence class in $Z_f$ for all $k$ and $i$. 
    It suffices to show that the elements of $Z_{f_k}$ and $Z_{f_{k'}}$ are not equivalent for $k \ne k'$. 
    Let $\beta \in Z_{f_k}$ and $\beta' \in Z_{f_{k'}}$, and  
    let $g_k$ and $g_{k'}$ be the minimal polynomials of $\beta$ and $\beta'$, respectively. 
    By \cref{thm:theorem of product}, the $\varphi$-polygons of $g_k$ and $g_{k'}$ have the same slope as $S$, 
    and $(g_k)_{S_{k}'} = \psi_k^{b_k}$ and $(g_{k'})_{S_{k'}'} = \psi_{k'}^{b_{k'}}$ for some $b_k$ and $b_{k'}$, where $S_{k}'$ and $S_{k'}'$ are the $\varphi$-polygons of $g_k$ and $g_{k'}$, respectively. 
    Let $\alpha_k$ and $\alpha_{k'}$ be roots of $g_k^0$ and $g_{k'}^0$, respectively. 
    Then, $\overline{\pi^{-h}{\alpha_k}^e}, \overline{\pi^{-h}{\alpha_{k'}}^e} \in \field{F}_{\hat{T}}$ are roots of $\psi_k$ and $\psi_{k'}$, respectively. 
    Since $\psi_k$ and $\psi_{k'}$ are coprime, we have $\overline{\pi^{-h}{\alpha_k}^e} \ne \overline{\pi^{-h}{\alpha_{k'}}^e}$ in $\field{F}_{\hat{T}}$, and 
    \[
      v_K(\pi^{-h}{\alpha_k}^e - \pi^{-h}{\alpha_{k'}}^e) = 0. 
    \]
    This implies that 
    \[
      v_K(\pi^{-\frac{h}{e}}\alpha_k - \pi^{-\frac{h}{e}}\alpha_{k'}) = 0, 
    \]
    and we have 
    \[
      v_K(\alpha_k - \alpha_{k'}) = \frac{h}{e}. 
    \]
    By \cref{lem:valuation of resultant}, 
    \[
      v_K(\mathrm{Res}(g_k,g_k^0)) > b_k e \cdot b_k h. 
    \]
    Let $L_0$ be the compositum of the splitting fields of $g_k$ and $g_k^0$ over $\hat{T}$. 
    Since the Galois group $\Gal(L_0/\hat{T})$ acts transitively on $Z_{g_k}$,
    for any two roots $\beta_0, \beta_1 \in Z_{g_k}$, there exists $\sigma_0 \in \Gal(L_0/\hat{T})$ such that 
    $\sigma_0(\beta_0) = \beta_1$, and we have $v_K(\prod_{\alpha \in Z_{g_k^0}} (\beta_0 - \alpha)) = v_K(\sigma_0(\prod_{\alpha \in Z_{g_k^0}} (\beta_0 - \alpha))) = v_K(\prod_{\alpha \in Z_{g_k^0}} (\beta_1 - \alpha))$. 
    Hence, 
    \[
      v_K(\prod_{\alpha \in Z_{g_k^0}} (\beta - \alpha)) > b_k h. 
    \]
    Moreover, by \cref{thm:theorem of the polygon}, for any $\alpha \in Z_{g_k^0}$, we have 
    \begin{align}
      v_K(\beta - \alpha) &\ge \min\{v_K(\beta),v_K(\alpha)\} \\
      &=\frac{h}{e}. 
    \end{align}
    These inequalities show that there exists a root $\alpha_k$ of $g_k^0$ such that 
    \[
      v_K(\beta - \alpha_k) > \frac{h}{e}. 
    \]
    On the other hand, since
    \[
      \frac{h}{e} = v_K(\alpha_k - \alpha_{k'}) = v_K(\alpha_k - \beta + \beta - \alpha_{k'}), 
    \]
    we have 
    \[
      v_K(\beta - \alpha_{k'}) = \frac{h}{e}. 
    \]
    Similarly, the same holds for $\beta'$. Therefore, we have 
    \[
      v_K(\beta - \beta') = v_K(\beta - \alpha_k + \alpha_k - \beta') = \frac{h}{e}. 
    \]

    Therefore, we may assume that $f_S = \psi^a$ for some linear polynomial $\psi$. 
    In this case, $f^0(X) = (X^e - b)^a \quad (b \in \mathcal{O}_K)$. 
    Let $\alpha$ be a root of $f^0$, and $\zeta$ a primitive $e$-th root of unity. 
    Then, the set of roots of $f^0$ is $Z_{f^0} = \{\alpha,\alpha\zeta, \dots ,\alpha\zeta^{e-1}\}$. 
    By \cref{lem:valuation of resultant}, for $\beta \in Z_f$, there exists an $i \in \{0, \dots ,e-1\}$ such that 
    \[
      v_K(\beta - \alpha\zeta^i) > \frac{h}{e}. 
    \]
    Moreover, since $\mathfrak{p} \nmid e$, we have 
    \[
      v_K(\prod_{i \ne 0} (\alpha - \alpha\zeta^i)) = v_K(\alpha^{e - 1}\prod_{i \ne 0} (1 - \zeta^i)) = (e - 1) \cdot \frac{h}{e}. 
    \]
    For $i \ne 0$, we have 
    \[
      v_K(\alpha - \alpha\zeta^i) \ge \min\{v_K(\alpha),v_K(\alpha\zeta^i)\} = \frac{h}{e}. 
    \]
    Hence, 
    \[
      v_K(\alpha - \alpha\zeta^i) = \frac{h}{e} \label{formula:lem-1}. 
    \]

    Put $Z_f^i = \{\beta \in Z_f \mid v_K(\beta - \alpha\zeta^i) > \frac{h}{e}\} \quad (i = 0, \dots ,e - 1)$. 
    We will show that each $Z_f^i$ forms an equivalence class in $Z_f$. 
    For $\beta, \beta' \in Z_f^i$, we have 
    \begin{align}
      v_K(\beta - \beta') &= v_K(\beta - \alpha\zeta^i + \alpha\zeta^i - \beta') \\
      &= \min\{v_K(\beta - \alpha\zeta^i),v_K(\beta' - \alpha\zeta^i)\} \\
      &> \frac{h}{e}. 
    \end{align}
    Hence, $\beta$ and $\beta'$ are equivalent. 
    Let $i \ne j$, and let $\beta \in Z_f^i$ and $\beta' \in Z_f^j$. 
    Then, we have 
    \begin{align}
      v_K(\beta - \beta') &= v_K(\beta - \alpha\zeta^j + \alpha\zeta^j - \beta') \\
      &\ge \min\{v_K(\beta - \alpha\zeta^j),v_K(\beta' - \alpha\zeta^j)\}. 
    \end{align}
    By (\ref{formula:lem-1}), 
    \begin{align}
      v_K(\beta - \alpha\zeta^j) &= v_K(\beta - \alpha\zeta^i + \alpha\zeta^i - \alpha\zeta^j) \\
      &= \frac{h}{e}. 
    \end{align}
    This implies that 
    \[
      v_K(\beta - \beta') = \frac{h}{e}. 
    \]
    Therefore, $Z_f^i$ forms an equivalence class. 

    It remains to show that $\#Z^i_f = a$ for all $i$. 
    Let $T = \hat{T}(\zeta)$. Since $\mathfrak{p} \nmid e$, $T/K$ is unramified. 
    If $e = 1$, there is only one equivalence class, and the claim is clear. 
    Assume that $e \ne 1$. 
    Let $\beta_i \in Z_f^i$, and let $L$ and $L_0$ be the splitting fields of $f$ and $f^0$ over $T$, respectively. 
    Since $v_K(\alpha) = \frac{h}{e}$ and $e$ and $h$ are coprime, we have $v_K(\alpha^m) \notin \Z$ for each $m = 1, \dots ,e-1$, hence $\alpha^m \notin T$. 
    By Kummer theory, $f^0$ is irreducible over $T$, and for any $i$ and $j$, there exists $\sigma_{ij} \in \Gal(LL_0/T)$ such that $\sigma_{ij}(\alpha \zeta^i) = \alpha \zeta^j$. 
    Then, we have 
    \[
      v_K(\beta_i - \alpha\zeta^i) = v_K(\sigma_{ij}(\beta_i) - \alpha\zeta^j), 
    \]
    hence $\sigma_{ij}(\beta_i) \in Z^j_f$. This implies that $\#Z^j_f \ge \#Z^i_f$ for any $i$ and $j$, and so $\#Z^j_f = \#Z^i_f$. 
    Therefore, we have shown that $\#Z^i_f = a$. 
  \end{proof}
  
  \begin{remark}
    In the proof of this lemma, the assumption $i_K(f) = i_{\varphi}(f)$ is not required. 
  \end{remark}

  \begin{proof}[Proof of \cref{prop:claim 4}]
    Put $n = \deg \psi$. Then, the length, height, and degree of $S$ are given by 
    $E = \deg f = \deg \tilde{f} = nea, H = nah, d = na$. 

    First, we show that it is sufficient to consider the case that $n = 1$. 
    Let $\hat{T}$ be the unramified extension of $K$ of degree $n$. 
    Then, $\psi = \prod_{i = 1}^n \psi_i$ for some distinct linear polynomials $\psi_i \quad (i = 1, \dots ,n)$. 
    Let $\theta$ and $\beta$ be roots of $f$ and $\tilde{f}$, respectively. 
    By \cref{thm:theorem of Montes Nart}, $\hat{T} \subseteq K(\theta), \hat{T} \subseteq K(\beta)$. 
    Hence, we have $K(\theta) = \hat{T}(\theta), K(\beta) = \hat{T}(\beta)$. 
    By \cref{prop:prop of Montes Nart}, $\tilde{\psi}(\pi^{-h}\theta^e)$ is a uniformizer of $K(\theta)$, where $\tilde{\psi}$ is a lift of $\psi$. 
    Let $g \in \mathcal{O}_{\hat{T}}[X]$ be the minimal polynomial of $\theta$ over $\hat{T}$. 
    By \cref{thm:theorem of Ore}, $g_S = \psi_i^l$ for some $i$ and $l$. 
    Assume that $i = 1$. 
    Since $g^0(\theta) + g^1(\theta) = g(\theta) = 0$, we have 
    \[
      v_K(g^0(\theta)) = v_K(g^1(\theta)) \ge H + \frac{1}{e}. 
    \]
    Moreover, let $\tilde{\psi}_1$ be a lift of $\psi_1$, since $g^0(\theta) = \pi^H\tilde{\psi}_1(\pi^{-h}\theta^e)^l$, 
    we have 
    \[
      v_K(\tilde{\psi}_1(\pi^{-h}\theta^e)) \ge \frac{1}{le} \ge \frac{1}{ea}. 
    \]
    On the other hand, we have 
    \[
      v_K(\tilde{\psi}_1(\pi^{-h}\theta^e)) \le v_K(\tilde{\psi}(\pi^{-h}\theta^e)) \le \frac{1}{ea}. 
    \]
    It follows that all the inequalities become equalities, $l = a$, and $\tilde{\psi}_1(\pi^{-h}\theta^e)$ is a uniformizer of $\hat{T}(\theta)$. 
    By \cref{prop:prop of Montes Nart}, the equality $i_{\hat{T}}(g) = i_{\varphi}(g)$ holds. 
    Therefore, it is sufficient to prove the proposition for $\hat{T}$ and $g$. 

    Hence, we can suppose that $\deg \psi = 1$. Let $\theta$ be a root of $f$. 
    Let $f(X) - \tilde{f}(X) = \sum_{i = 0}^E c_{E - i}X^i$. 
    Since $f_S = \tilde{f}_S, f_{\tilde{S}}^1 = \tilde{f}_{\tilde{S}}^1$, 
    for each $i$, we have 
    \[
      v_K(c_{E - i}) \ge H - \frac{h}{e} \cdot i + \frac{2}{e}. 
    \]
    Thus, we have 
    \[
      v_K(\tilde{f}(\theta)) = v_K(f(\theta) - \tilde{f}(\theta)) \ge H + \frac{2}{e} > H + \frac{1}{e}. 
    \]
    By \cref{thm:theorem of Montes Nart}, the ramification index of $K(\theta)/K$ is equal to $ea$. 
    Denote the set of roots of $\tilde{f}$ by $\{\beta_j \mid j = 1, \dots ,ea\}$. 
    For all $j$, we have 
    \[
      v_K(\theta - \beta_j) \ge \min\{v_K(\theta),v_K(\beta_j)\} = \frac{h}{e}. 
    \]
    If $v_K(\theta - \beta_{j_1}) > \frac{h}{e}$ and $v_K(\beta_{j_1} - \beta_{j_2}) = \frac{h}{e}$, we have 
    \[
      v_K(\theta - \beta_{j_2}) = v_K(\theta - \beta_{j_1} + \beta_{j_1} - \beta_{j_2}) = \frac{h}{e}. 
    \]
    By \cref{lem:classification of roots}, there are at most $a$ roots such that $v_K(\theta - \beta_j) \ge \frac{h}{e} + \frac{1}{ea}$. 
    Noting that $\tilde{f}$ has $ea$ roots, we obtain from above inequalities that there exists $j_0 \in \{1, \dots E\}$ such that 
    \[
      v_K(\theta - \beta_{j_0}) > \frac{h}{e} + \frac{1}{ea} \label{formula:prop4-2}. 
    \]
    Assume that $j_0 = 1$, and denote $\beta = \beta_1$. 

    We consider the derivative $\tilde{f}'$ of $\tilde{f}$. 
    Since $\tilde{f}^0(X) = \pi^H\tilde{\psi}(\pi^{-h}X^e)^a$, we have 
    \[
      (\tilde{f}^0)'(X) = ea\pi^{(a-1)h}X^{e - 1}\tilde{\psi}(\pi^{-h}X^e)^{a - 1}. 
    \]
    Noting that $\mathfrak{p} \nmid ea$ and $\tilde{\psi}(\pi^{-h}\beta^e)$ is a uniformizer of $K(\beta)$, it follows that 
    \[
      v_K((\tilde{f}^0)'(\beta)) = ah - \frac{h}{e} +\frac{a - 1}{ea}. 
    \]
    Since the points of  the $\varphi$-polygon of $\tilde{f}^1$ lie on $\tilde{S}$, we have 
    \[
      v_K((\tilde{f}^1)'(\beta)) \ge ah - \frac{h}{e} + \frac{1}{e} > v_K((\tilde{f}^0)'(\beta)). 
    \]
    Hence, 
    \begin{align}
      \sum_{i \ne 1} v_K(\beta - \beta_i) &= v_K(\tilde{f}'(\beta)) \\ 
      &=v_K((\tilde{f}^0)'(\beta)) \\
      &= ah - \frac{h}{e} + \frac{a - 1}{ea}
    \end{align}
    Since $v_K(\beta - \beta_i) \ge \frac{h}{e}$ for $i \ne 1$, we have 
    \[
      v_K(\beta - \beta_i) \le \frac{h}{e} + \frac{a - 1}{ea}. 
    \]
    By \cref{lem:classification of roots}, there are just $a - 1$ roots $\beta_i$ satisfying $v_K(\beta - \beta_i) > \frac{h}{e}$. 
    Consequently, for all $i = 2, \dots ,ea$, 
    \[
      v_K(\beta - \beta_i) \le \frac{h}{e} + \frac{1}{ea}. \label{formula:prop4-3}
    \]

    From (\ref{formula:prop4-2}) and (\ref{formula:prop4-3}), for all $i \ne 1$, 
    we have 
    \[
      v_K(\theta - \beta) > v_K(\beta - \beta_i). 
    \]
    By Krasner's lemma, this gives $K(\beta) \subset K(\theta)$. 
    We have equality since $f$ and $\tilde{f}$ are irreducible over $K$ of the same degree. 

    It remains to show the uniqueness of $\beta$. 
    If there exists another root $\beta'$ of $\tilde{f}$ satisfying (\ref{formula:prop4-2}), then we have 
    \[
      v_K(\beta - \beta') > \frac{h}{e} + \frac{1}{ea}. 
    \]
    This contradicts (\ref{formula:prop4-3}). 
  \end{proof}


  \begin{proof}[Proof of \cref{thm:our main theorem}]
  We now prove \cref{thm:our main theorem}. 
  We use the notation of \cref{sec:Kölle Schmid}. 

  By \cref{prop:claim 2}, $\tilde{f}$ is separable. 
  By \cref{thm:theorem of the polygon}, the $\mathfrak{p}$-adic valuations of all $de$ roots of $f$ are equal to $m$. 

  Assume that $\mathfrak{p} \nmid ea$. 
  (i) follows from \cref{prop:claim 3} and \cref{prop:claim 4}. 
  By (i), the splitting fields of $f$ and $\tilde{f}$ coincide, say $L$.
  Hence, $G_{\mathfrak{P}}$ is isomorphic to $\Gal_{K_{\mathfrak{p}}}(\tilde{f})$ as a finite group. 
  Let $\theta$ be a root of $f$, and $g(X) \in K_{\mathfrak{p}}[X]$ the minimal polynomial of $\theta$ over $K_{\mathfrak{p}}$. 
  Let $\tilde{g}$ be the irreducible factor of $\tilde{f}$ corresponding to $g$ under the bijection of \cref{prop:claim 3}, and 
  $\beta$ be the root of $\tilde{g}$ corresponding to $\theta$ under the bijection of \cref{prop:claim 4}. 
  Since $\deg g = \deg \tilde{g}$, the $G_{\mathfrak{P}}$-orbit of $\theta$ has the same length as the $\Gal_{K_{\mathfrak{p}}}(\tilde{f})$-orbit of $\beta$. 
  Moreover, the stabilizers of $\theta$ and $\beta$ can be written as $\Gal(L/K_{\mathfrak{p}}(\theta))$ and $\Gal(L/K_{\mathfrak{p}}(\beta))$, respectively. 
  Since $K_{\mathfrak{p}}(\theta) = K_{\mathfrak{p}}(\beta)$, it follows that the stabilizers of $\theta$ and $\beta$ coincide. 
  These imply (ii). 

  Put $G' = \Gal(L/K_{\mathfrak{p}})$. By (ii), we may identify $G_{\mathfrak{P}}$ with $G'$, together with their action on $Z_f$ (or $Z_{\tilde{f}}$). 
  Let $\hat{T}$ be the maximal unramified subextension of $L/K_{\mathfrak{p}}$, then $I_{\mathfrak{P}} = \Gal(L/\hat{T})$. 

  By \cref{thm:theorem of Montes Nart}, the factorization of \cref{thm:theorem of Ore} 
  \[
    \tilde{f}(X) = \tilde{f}_1(X) \cdots \tilde{f}_r(X)
  \]
  is an irreducible factorization. Hence, $Z_f$ has $r$ orbits of lengths $d_1 ea_1, \dots ,d_r ea_r$ under the action of $G$. 
  Moreover, since $L/\hat{T}$ is tamely ramified, $G_{\mathfrak{P}}$ is a metacyclic group. 
  Let $\tau$ be a generator of $I_{\mathfrak{P}}$, and $\sigma$ be a lift of a generator of $G_{\mathfrak{P}}/I_{\mathfrak{P}}$. 
  Then we have $\sigma^{-1}\tau\sigma = \tau^{N\mathfrak{p}}$. 

  Let $\varepsilon \in \bar{K}_{\mathfrak{p}}$ be a primitive $ea$-th root of unity. 
  Then $[K_{\mathfrak{p}}(\varepsilon):K_{\mathfrak{p}}] = \omega$. 
  Let $\beta$ be a root of $\tilde{f}$, and $T'$ be the maximal unramified subextension of $K_{\mathfrak{p}}(\beta)/K_{\mathfrak{p}}$. 
  Since $K_{\mathfrak{p}}(\beta)/K_{\mathfrak{p}}$ is tamely ramified and by the theorem of Pauli and Roblot (\cite[Theorem 7.2]{MR1836924}), 
  there exists an element $u \in T'$ such that $K_{\mathfrak{p}}(\beta)$ coincides with the extension over $T'$ defined by the polynomial $X^{ea_i} - u$. 
  Let $\alpha$ be a root of $X^{ea_i} - u$. Then the set of roots of $X^{ea_i} - u$ can be written as $\{\alpha,\alpha\varepsilon_i,\dots,\alpha\varepsilon_i^{ea_i-1}\}$, 
  where $\varepsilon_i$ is a primitive $ea_i$-th root of unity. 
  Since the normal closure of $K_{\mathfrak{p}}(\beta)$ is contained in $L$, $\varepsilon_i \in L$ for all $i=1, \dots ,r$. 
  Hence, $\varepsilon \in L$. 
  Since $\mathfrak{p} \nmid ea$, $K_{\mathfrak{p}}(\varepsilon)/K_{\mathfrak{p}}$ is unramified. Hence, $K_{\mathfrak{p}}(\varepsilon) \subset \hat{T}$. 
  It follows that $\omega$ divides $|G_{\mathfrak{P}}/I_{\mathfrak{P}}|$. 
  Moreover, let $\beta_i$ be a root of $\tilde{f}_i$. 
  By \cref{thm:theorem of Montes Nart}, the degree over $K_{\mathfrak{p}}$ of the maximal unramified subextension of $K_{\mathfrak{p}}(\beta_i)/K_{\mathfrak{p}}$ 
  is equal to $d_i$. It follows that $d_i$ divides $|G_{\mathfrak{P}}/I_{\mathfrak{P}}|$ for each $i$. 
  Therefore $\mu$ divides $|G_{\mathfrak{P}}/I_{\mathfrak{P}}|$. 

  Since $d_i$ divides $|G_{\mathfrak{P}}/I_{\mathfrak{P}}|$ for each $i$, $\tilde{f}_S$ factors as a product of linear polynomials over $\field{F}_{\hat{T}}$. 
  Since $\hat{T}/K_{\mathfrak{p}}$ is unramified, the $\varphi$-polygon of $\tilde{f}$ over $\hat{T}$ coincides with that over $K_{\mathfrak{p}}$. 
  This implies that the equality $i_{\hat{T}}(\tilde{f}) = i_{\varphi}(\tilde{f})$ also holds. 
  Indeed, the polynomial $\tilde{f}^1$ defined over $\hat{T}$ coincides with that defined over $K_{\mathfrak{p}}$. 
  Noting that $\tilde{f}^1_{\tilde{S}}$ is a polynomial over $\field{F}_{K_{\mathfrak{p}}}$, it follows from \cref{thm:equivalent condition of Montes Nart} that 
  $i_{\hat{T}}(\tilde{f}) = i_{\varphi}(\tilde{f})$. 
  By \cref{thm:theorem of Montes Nart}, $\tilde{f}$ factors as a product of irreducible polynomials of degree $ea_i$. 
  Thus, $I_{\mathfrak{P}}$-orbits of $Z_{\tilde{f}}$ have lengths $ea_i$. 
  Since $L/K_{\mathfrak{p}}$ is tamely ramified, $I_{\mathfrak{P}}$ is a cyclic group. 
  These imply (iii). 

  There exists a unique subfield $T$ of $\hat{T}$ such that $[T:K_{\mathfrak{p}}] = \mu$. 
  Let $\beta$ be a root of $\tilde{f}$, and suppose that $\tilde{f}_j$ is the minimal polynomial of $\beta$. 
  Since $d_j \mid \mu$, $(\tilde{f}_j)_S$ factors as a product of linear polynomials over $\field{F}_T$. 
  By \cref{thm:theorem of Montes Nart}, $T(\beta)/T$ is a totally ramified extension of degree $ea_j$. 
  Since $\varepsilon \in T$, $T(\beta)/T$ is a Galois extension by the theorem of Pauli and Roblot and Kummer theory. 
  On the other hand, $\hat{T}/T$ is unramified, and hence $T(\beta)$ and $\hat{T}$ are linearly disjoint over $T$. 
  It follows that $T(\beta)\hat{T}/T(\beta)$ is a cyclic extension of degree $[\hat{T}:T]$. 
  Noting that $L$ is the compositum of all these $T(\beta)$ and $T(\beta)/T$ is a Galois extension of degree $ea_i$, 
  the exponent of its Galois group $\Gal(L/T)$ is $ea$. 
  Hence, the exponent of $\Gal(T(\beta)\hat{T}/T)$ divides $ea$. 
  Since $\Gal(T(\beta)\hat{T}/T(\beta))$ is a cyclic subgroup of $\Gal(T(\beta)\hat{T}/T)$, 
  its degree $[\hat{T}:T]$ divides $ea$. Therefore $|G_{\mathfrak{P}}/I_{\mathfrak{P}}|$ divides $\mu ea$. 
  Thus, (iv) has been proved. 

  \end{proof}


\begin{thebibliography}{99}
\bibitem[MR1152908]{MR1152908} \href{https://doi.org/10.1016/0021-8693(92)90071-S} {J.~Montes and E.~Nart. \newblock On a theorem of {O}re. \newblock {\em J. Algebra}, 146(2):318--334, 1992.}
\bibitem[MR1836924]{MR1836924} \href{https://doi.org/10.1090/S0025-5718-01-01306-0} {S.~Pauli and X.-F.~c. Roblot. \newblock On the computation of all extensions of a {$p$}-adic field of a given degree. \newblock {\em Math. Comp.}, 70(236):1641--1659, 2001.}
\bibitem[MR2102807]{MR2102807} \href{https://doi.org/10.4064/aa115-1-6} {M.~K\"olle and P.~Schmid. \newblock Computing {G}alois groups by means of {N}ewton polygons. \newblock {\em Acta Arith.}, 115(1):71--84, 2004.}
\bibitem[MR2557856]{MR2557856} \href{https://doi.org/10.4064/aa140-1-7} {M.~K\"olle and P.~Schmid. \newblock Correction to ``{C}omputing {G}alois groups by means of {N}ewton polygons'' [mr2102807]. \newblock {\em Acta Arith.}, 140(1):101--103, 2009.}
\bibitem[index]{index} N.~Suzuki. \newblock Calculating the index of equation orders using Newton polygons. \newblock {\em Kodai Math. J.}, in press.
\end{thebibliography}
\end{document}
